This work addressed the mechanical reengineering and automation of an existing Cartesian robotic arm for its application in the packaging line of the Process Design Laboratory at Universidad del Valle de Guatemala. The existing system presented mechanical limitations along the X, Y, and Z axes, including high friction, misalignment, and irregular motion, which affected its accuracy, repeatability, and ability to operate stably. To address these deficiencies, the structural components and translation systems were partially redesigned using computer-aided design tools. Guide shafts and linear bearings were also incorporated to improve the smoothness and stability of motion.

The automation of the system was developed through an embedded control architecture responsible for coordinating the motors, sensors, communication devices, and peripherals. Limit switches were integrated to perform automatic calibration and establish motion references, while a rotary encoder was incorporated to measure the speed of the conveyor belt and facilitate process synchronization. In addition, a computer vision system was implemented to identify objects with previously defined characteristics, such as shape and color, and to determine the information required to perform detection, classification, manipulation, and transfer tasks through a pick-and-place sequence. The robotic arm also included a manual operating mode using a wireless controller and a local interface to display the calibration status, operating mode, and actions performed by the system.

The methodology was organized into phases consisting of evaluation of the existing mechanism, mechanical redesign, manufacturing and assembly, development of the electronic architecture, motion-control programming, integration of the computer vision system, and experimental validation. The tests were conducted in a controlled academic environment and allowed variables such as positioning accuracy, repeatability, response time, object-detection capability, and overall system performance to be evaluated. For this purpose, repetitive motions and operating cycles were performed, the achieved positions were compared with the target positions, and the variation among different executions was analyzed.

As a result, a Cartesian robotic arm with smoother, more stable, and more repeatable motion was obtained, capable of performing automated product manipulation and transfer tasks within a simulated packaging line. Likewise, the system was developed as an academic platform for teaching and experimentation in areas related to industrial automation, robotics, embedded systems, sensors, and computer vision. Although the project was not validated under continuous industrial production conditions, it demonstrated the feasibility of the implemented mechanical improvements and automation within a laboratory environment.